\documentclass[12pt]{article}
\usepackage[utf8]{inputenc}
\usepackage[T1]{fontenc}
\usepackage{amsmath, amssymb}
\usepackage{hyperref}
\usepackage{geometry}
\geometry{margin=1in}

\title{Sample Document for LaTeX Commenter}
\author{Amin Aghababaei}
\date{\today}

\usepackage[colorinlistoftodos]{todonotes}
\begin{document}

\maketitle

\begin{abstract}
This document is a sandbox for testing the LaTeX Commenter extension.
Select any sentence or phrase, then use to annotate it. Hover over highlighted text to see comments inline.
\end{abstract}

% ─────────────────────────────────────────────────────────────────────────────
\section{Introduction}
% ─────────────────────────────────────────────────────────────────────────────

Machine learning has become one of the most transformative technologies of the
past decade, enabling applications that were previously considered impossible.
From image recognition to natural language processing, neural networks now
underpin a broad spectrum of real-world systems.

Despite these advances, a fundamental challenge remains: models trained on one
distribution often fail catastrophically when deployed in a different
environment. This phenomenon, known as distribution shift, is a critical
bottleneck in deploying reliable AI systems at scale.

In this paper, we propose a novel regularization scheme that improves
out-of-distribution generalization without requiring access to target-domain
data at training time. Our method achieves state-of-the-art results on five
standard benchmarks while adding negligible computational overhead.

% ─────────────────────────────────────────────────────────────────────────────
\section{Related Work}
% ─────────────────────────────────────────────────────────────────────────────

Early approaches to domain adaptation relied on explicit feature alignment
between source and target domains~\cite{pan2010survey}. These methods typically
minimized the Maximum Mean Discrepancy (MMD) between feature distributions,
which required paired or unpaired samples from both domains.

More recent work has explored invariant risk minimization (IRM), which seeks
representations that achieve the same optimal classifier across all training
environments. While theoretically elegant, IRM is notoriously difficult to
optimize in practice and sensitive to the choice of penalty coefficient.

Contrastive learning methods have also shown promise for learning
distribution-robust representations. By pulling together augmented views of
the same sample and pushing apart views from different samples, these
approaches encourage features that are invariant to nuisance transformations.

% ─────────────────────────────────────────────────────────────────────────────
\section{Method}
% ─────────────────────────────────────────────────────────────────────────────

\subsection{Problem Setup}

Let $\mathcal{D} = \{(x_i, y_i)\}_{i=1}^{N}$ denote a training dataset drawn
from source distribution $P_s$. At test time, inputs are drawn from a
potentially different distribution $P_t$. Our goal is to learn a hypothesis
$h: \mathcal{X} \to \mathcal{Y}$ that minimizes the expected loss under $P_t$,
given access only to samples from $P_s$.

\subsection{Proposed Regularizer}

We introduce a \emph{diversity penalty} $\Omega(\theta)$ that encourages the
model to rely on a diverse set of input features rather than a single
high-signal shortcut. Formally,

\begin{equation}
  \mathcal{L}(\theta) = \mathbb{E}_{(x,y)\sim P_s}\bigl[\ell(h_\theta(x), y)\bigr]
  + \lambda\,\Omega(\theta),
\end{equation}

where $\lambda > 0$ is a scalar hyperparameter and $\Omega$ is defined as the
negative entropy of the normalized gradient magnitudes with respect to input
features.

\subsection{Optimization}

We train using stochastic gradient descent with momentum. The penalty term
$\Omega$ is computed over mini-batches and back-propagated through the feature
extractor. We find that a warm-up schedule for $\lambda$ is essential for
stable training, particularly in early epochs when the model has not yet
learned meaningful representations.

% ─────────────────────────────────────────────────────────────────────────────
\section{Experiments}
% ─────────────────────────────────────────────────────────────────────────────

We evaluate our method on the DomainBed benchmark suite, which includes
five datasets: PACS, VLCS, OfficeHome, TerraIncognita, and DomainNet.
All experiments use a ResNet-50 backbone pretrained on ImageNet.

\begin{table}[h]
  \centering
  \caption{Average out-of-distribution accuracy (\%) across five datasets.}
  \begin{tabular}{lcc}
    \hline
    \textbf{Method}  & \textbf{Avg.\ Accuracy} & \textbf{Std.\ Dev.} \\
    \hline
    ERM              & 64.3                    & 0.6 \\
    IRM              & 63.5                    & 0.8 \\
    GroupDRO         & 65.1                    & 0.7 \\
    \textbf{Ours}    & \textbf{67.4}           & \textbf{0.5} \\
    \hline
  \end{tabular}
\end{table}

Our method consistently outperforms all baselines. Notably, it shows the
largest gains on datasets with strong spurious correlations, which is consistent
with the motivation behind our diversity penalty.

% ─────────────────────────────────────────────────────────────────────────────
\section{Conclusion}
% ─────────────────────────────────────────────────────────────────────────────

We presented a simple and effective regularization technique for improving
out-of-distribution generalization. Our experiments demonstrate consistent
improvements across diverse benchmarks. Future work will explore extending
this approach to sequence models and large language models, where distribution
shift is equally prevalent but less well-studied.

\bibliographystyle{plain}
\begin{thebibliography}{9}
\bibitem{pan2010survey}
  Sinno Jialin Pan and Qiang Yang.
  \textit{A Survey on Transfer Learning}.
  IEEE Transactions on Knowledge and Data Engineering, 22(10):1345--1359, 2010.
\end{thebibliography}

\end{document}
